\documentclass[11pt]{article}
\usepackage[margin=1in]{geometry}
\usepackage{amsmath,amssymb,amsthm,mathtools,bm}
\usepackage{enumitem}
\usepackage[dvipsnames]{xcolor}
\usepackage{hyperref}
\hypersetup{colorlinks=true,linkcolor=NavyBlue,citecolor=NavyBlue,urlcolor=NavyBlue}

\theoremstyle{plain}
\newtheorem{theorem}{Theorem}[section]
\newtheorem{proposition}[theorem]{Proposition}
\newtheorem{lemma}[theorem]{Lemma}
\newtheorem{corollary}[theorem]{Corollary}
\theoremstyle{definition}
\newtheorem{definition}[theorem]{Definition}
\newtheorem{assumption}[theorem]{Assumption}
\theoremstyle{remark}
\newtheorem{remark}[theorem]{Remark}

% macros
\newcommand{\Tr}{\operatorname{Tr}}
\newcommand{\id}{\mathbb{1}}
\newcommand{\E}{\mathbb{E}}
\newcommand{\Var}{\operatorname{Var}}
\newcommand{\ket}[1]{\lvert #1 \rangle}
\newcommand{\bra}[1]{\langle #1 \rvert}
\newcommand{\braket}[2]{\langle #1 \vert #2 \rangle}
\newcommand{\proj}[1]{\lvert #1 \rangle\!\langle #1 \rvert}
\newcommand{\SWAP}{\mathbb{S}}
\newcommand{\Hs}{\mathcal{H}}
\newcommand{\FQ}{F_{Q}}
\newcommand{\norm}[1]{\lVert #1 \rVert}
\newcommand{\eps}{\varepsilon}
\newcommand{\code}[1]{\texttt{#1}}
\DeclareMathOperator{\poly}{poly}
\newcommand{\bigO}{\mathcal{O}}

\title{\textbf{Information-Theoretic Advantage Guarantees for a Sewn,\\
Instantaneously-Deep Quantum-Memory Reservoir}\\[4pt]
\large Formal separations on five tasks, with explicit Willow-platform constants}
\author{ID--QND--RC project --- theory companion}
\date{}

\begin{document}
\maketitle

\begin{abstract}
We give quantum-information-theoretic guarantees of advantage for the five applications implemented on the sewn, instantaneously-deep quantum-memory reservoir (the ID--QND--RC platform) running on Google Willow. For each task we (i) fix a precise computational model of the platform, (ii) prove an \emph{achievability} bound for the platform's coherent two-copy / entangled / recurrent read-out, (iii) state and prove (or cite, with proof sketch) a matching \emph{lower bound} for every strategy lacking quantum memory, and (iv) record an explicit \emph{separation corollary} with the platform's parameters $(N,|A|,M,\eta)$ and a noise-robustness condition under which the separation survives Willow-class error rates. Tasks~1 and~3 (estimating non-linear functionals of states; tracking entanglement) enjoy \emph{unconditional, exponential} sample separations rooted in the with-/without-quantum-memory dichotomy. Task~2 (sensing) enjoys an \emph{unconditional, quadratic} metrological separation via the quantum Fisher information. Task~4 (decoding) enjoys an \emph{unconditional latency / learnability} guarantee (not an exponential data separation; the inputs are classical). Task~5 (generative modelling) enjoys an \emph{unconditional memory} separation and a \emph{conditional} (complexity-theoretic) sampling-hardness separation. Throughout we are explicit about which claims are theorems proved here, which are theorems cited from the literature, and which are conditional.
\end{abstract}

\tableofcontents

%==========================================================================
\section{The platform, formally}\label{sec:platform}

\subsection{State space, dynamics, edge of chaos}
The reservoir occupies a system register $S$ of $N$ qubits, $\Hs_S=(\mathbb C^2)^{\otimes N}$, $d:=2^N$, together with an ancilla register used for read-out. The reservoir map is generated by the \emph{critical-phase} step
\begin{equation}
U_{g}\;=\;\Big(\prod_{i} R^{x}_i(\theta_x)\Big)\Big(\prod_{i} R^{z}_i(\theta_z)\Big)\Big(\prod_{\langle i,j\rangle}\mathrm{CZ}_{ij}^{\,2g/\pi}\Big),
\qquad g=g^\star=0.30,
\label{eq:Ustep}
\end{equation}
tuned to the edge-of-chaos critical point $g^\star$ where the operator entanglement grows fastest while remaining trainable. All gates in \eqref{eq:Ustep} are realised as genuine \code{cirq} operations and transpile to the Willow native set $\{\mathrm{CZ},\mathrm{PhasedXZ}\}$.

\subsection{The sewn collapse-free read-out}
The defining capability of the platform is a read-out that extracts a local observable \emph{without collapsing} the memory. We formalise it through the deferred-measurement identity.

\begin{lemma}[Coherent-copy / collapse-free read-out]\label{lem:sew}
Let $b\in\{X,Y,Z\}$ and let $\{\ket{e^b_0},\ket{e^b_1}\}$ be the eigenbasis of the Pauli $b$ on system qubit $j$. Couple a fresh ancilla $a$ (in $\ket 0$) by the basis-change $V_b$ that maps $b\to Z$, a $\mathrm{CNOT}_{j\to a}$, and $V_b^\dagger$; then measure $a$ in the computational basis and discard the outcome. The induced channel on $\Hs_S$ is the dephasing channel in basis $b$ on qubit $j$,
\begin{equation}
\mathcal{D}^{(b)}_j(\rho)\;=\;\Pi^b_0\,\rho\,\Pi^b_0+\Pi^b_1\,\rho\,\Pi^b_1,
\qquad \Pi^b_s=\proj{e^b_s},
\label{eq:dephase}
\end{equation}
and the recorded statistics of $a$ reproduce $\Tr(\Pi^b_1\rho)$ exactly. In particular the post-read-out reduced state on the \emph{other} $N-1$ qubits is undisturbed, and the diagonal of $\rho$ in basis $b$ on qubit $j$ is preserved: the memory persists.
\end{lemma}

\begin{proof}
Writing $\rho$ in the $b$-eigenbasis of qubit $j$ as $\rho=\sum_{s,s'}\rho_{ss'}\proj{e^b_s}\!{}_{,j}\,(\cdots)$, the unitary $V_b^\dagger\,\mathrm{CNOT}_{j\to a}\,V_b$ acts as $\ket{e^b_s}_j\ket0_a\mapsto\ket{e^b_s}_j\ket{s}_a$. Hence the joint state becomes $\sum_{s,s'}\rho_{ss'}\ket{e^b_s}\!\bra{e^b_{s'}}_j\otimes\ket s\!\bra{s'}_a\otimes(\cdots)$. Measuring $a$ in the computational basis and tracing it out applies $\sum_s\proj{s}_a(\cdot)\proj{s}_a$, killing the $s\neq s'$ coherences on $a$ and therefore on $j$, which is exactly \eqref{eq:dephase}; the probability of outcome $s$ is $\rho_{ss}=\Tr(\Pi^b_s\rho)$. The map acts as identity on every qubit other than $j$. \end{proof}

\begin{remark}
Equation~\eqref{eq:dephase} is a \emph{unital, basis-diagonal} channel, not a projective collapse: repeated sewn read-outs in commuting bases leave a nontrivial quantum memory that continues to evolve under \eqref{eq:Ustep}. This is precisely the resource that removes the ``rewinding'' overhead of standard quantum reservoir computing, in which a projective read-out forces re-execution of the circuit to re-inject identical data.
\end{remark}

\subsection{Instantaneous depth}
\begin{assumption}[Instantaneous depth; \cite{LSQC,GenAdv}]\label{ass:idepth}
A target reservoir map of logical depth $D$ (a degree-$\le 2^{D}$ polynomial feature map) is realised at \emph{physical} circuit depth $\bigO(1)$ in $D$ by trading depth for ancilla width via local inversions / the sewing construction, at an ancilla cost polynomial in $D$ and in the light-cone size.
\end{assumption}
Assumption~\ref{ass:idepth} is what keeps every application circuit shallow enough to be Willow-native and to satisfy the noise budgets stated below.

\subsection{Two-copy access primitive}
Tasks~1 and~3 estimate functionals quadratic in $\rho$. The platform realises coherent two-copy access either by an ancilla-controlled $\mathrm{SWAP}$ (in exact simulation) or, on hardware, by a Bell-basis measurement (one $\mathrm{CNOT}+H$ per matched pair, all native).

\begin{lemma}[Two-copy purity primitive]\label{lem:swap}
Let $\SWAP$ be the swap on two $n$-qubit copies and $\rho$ an $n$-qubit state. The SWAP test with control ancilla $a$ outputs $a=0$ with probability $p_0=\tfrac12\big(1+\Tr(\rho^2)\big)$. Equivalently, a Bell-basis measurement on the $n$ matched qubit pairs yields a $\pm1$ variable $\xi$ with $\E[\xi]=\Tr(\rho^2)$. Both estimators are unbiased for $\Tr(\rho^2)$ and bounded in $[-1,1]$.
\end{lemma}
\begin{proof}
$\Tr(\SWAP\,\rho\otimes\rho)=\Tr(\rho^2)$. The Hadamard test on the unitary $\SWAP$ gives $p_0=\tfrac12(1+\Re\Tr(\SWAP\rho^{\otimes2}))=\tfrac12(1+\Tr\rho^2)$. For the Bell variant, $\SWAP=\sum_{k}(-1)^{[k\in\text{singlets}]}\Pi_k$ over the two-qubit Bell projectors per pair, so the product of single-pair $\pm1$ outcomes has expectation $\Tr(\SWAP\rho^{\otimes2})=\Tr\rho^2$.\end{proof}

\subsection{Platform parameters and the noise model}\label{sec:noise}
We use a Pauli/depolarising surrogate for Willow-class hardware: two-qubit gate error $\eta_2$, one-qubit error $\eta_1$, read-out error $\eta_{\mathrm{ro}}$, with representative values $\eta_2\approx 3\times10^{-3}$, $\eta_1\approx 10^{-3}$, $\eta_{\mathrm{ro}}\approx 7\times10^{-3}$, coherence $T_1,T_2\gtrsim 10^2~\mu$s, on a $105$-qubit chip \cite{Willow}. For a circuit with $g_2$ two-qubit gates, $g_1$ one-qubit gates and $r$ measured qubits, the diamond-distance deviation of the realised channel from the ideal obeys the subadditive bound
\begin{equation}
\Delta\;:=\;\tfrac12\norm{\widetilde{\mathcal E}-\mathcal E}_{\diamond}\;\le\;\eta_2 g_2+\eta_1 g_1+\eta_{\mathrm{ro}} r\;=:\;\Gamma .
\label{eq:errbudget}
\end{equation}
Any expectation value read off the device is therefore biased by at most $\Gamma$ (in the relevant normalisation), which is the single quantity that enters every noise-robustness condition below.

%==========================================================================
\section{Task 1: learning non-linear functionals from quantum experiments}\label{sec:task1}

\subsection{Model and the two access classes}
\begin{definition}[Access classes]\label{def:access}
A learner receives $T$ copies of an unknown $\rho\in\Hs_S$.
\emph{Incoherent (single-copy, ``without quantum memory'')}: each copy is measured by an arbitrary POVM, possibly chosen adaptively from past classical outcomes; the algorithm is any function of the classical transcript. This class contains \emph{all} classical-shadow protocols (Pauli or global Clifford).
\emph{Coherent two-copy (the platform)}: pairs of copies are jointly measured (e.g.\ by Lemma~\ref{lem:swap}); a persistent quantum memory (Lemma~\ref{lem:sew}) holds one copy while the next arrives.
\end{definition}

The canonical functional is the purity $P(\rho)=\Tr(\rho^2)$; testing $P=1$ versus $P\le\tfrac12$ already separates the classes.

\subsection{Achievability on the platform}
\begin{proposition}[Dimension-independent two-copy estimation]\label{prop:t1up}
Using $M$ two-copy measurements (Lemma~\ref{lem:swap}), the empirical mean $\widehat P=\frac1M\sum_{m}\xi_m$ is unbiased for $\Tr(\rho^2)$ and, for any $\eps,\delta\in(0,1)$,
\begin{equation}
\Pr\!\big[\,|\widehat P-\Tr\rho^2|\ge \eps\,\big]\le \delta
\qquad\text{whenever}\qquad
M\;\ge\; M_Q(\eps,\delta):=\Big\lceil \frac{\ln(2/\delta)}{2\eps^{2}}\Big\rceil,
\label{eq:t1up}
\end{equation}
\emph{independently of the number of qubits $N$}.
\end{proposition}
\begin{proof}
Unbiasedness is Lemma~\ref{lem:swap}. Each $\xi_m\in[-1,1]$ i.i.d.; Hoeffding's inequality gives $\Pr[|\widehat P-\E\xi|\ge\eps]\le 2\exp(-2M\eps^2/(b-a)^2)$ with $b-a=2$, i.e.\ $\le2\exp(-M\eps^2/2)$. Setting this $\le\delta$ yields \eqref{eq:t1up}. No dependence on $d=2^N$ enters.\end{proof}

\subsection{Lower bound without quantum memory}
We first isolate the mechanism: against the maximally mixed null, single-copy statistics carry \emph{no first-order} signal, so distinguishing requires resolving second moments, which costs $\sqrt d$ copies.

\begin{lemma}[Haar two-design moments]\label{lem:2design}
For Haar-random $\ket\psi\in\mathbb C^{d}$,
$\E\!\big[\proj\psi\big]=\id/d$ and
$\E\!\big[\proj\psi^{\otimes2}\big]=\dfrac{\id+\SWAP}{d(d+1)}.$
\end{lemma}
\begin{proof}
The first is Schur's lemma. For the second, $\E[\proj\psi^{\otimes2}]$ commutes with all $U\otimes U$, hence by Schur--Weyl it is a combination $\alpha\id+\beta\SWAP$ supported on the symmetric subspace; normalisation $\Tr=1$ and support on the symmetric projector $\tfrac12(\id+\SWAP)$ of dimension $\binom{d+1}{2}$ give $\E[\proj\psi^{\otimes2}]=\tfrac{\id+\SWAP}{d(d+1)}$.\end{proof}

\begin{lemma}[First-moment blindness]\label{lem:blind}
Fix any POVM $\{E_x\}$. Under $H_0:\rho=\id/d$ the outcome law is $p_0(x)=\Tr(E_x)/d$. Under $H_1:\rho=\proj\psi$ with $\psi$ Haar, $\E_\psi\,p_\psi(x)=p_0(x)$ for every $x$.
\end{lemma}
\begin{proof}
$\E_\psi\,p_\psi(x)=\Tr\!\big(E_x\,\E_\psi\proj\psi\big)=\Tr(E_x\,\id/d)=\Tr(E_x)/d=p_0(x)$ by Lemma~\ref{lem:2design}.\end{proof}

Because the means coincide, the $\chi^2$ between the $T$-copy laws is governed entirely by the (shared-$\psi$) second moment. Writing $Q=\E_\psi\,p_\psi^{\otimes T}$ and $P_0=p_0^{\otimes T}$,
\begin{align}
1+\chi^2(Q\Vert P_0)
&=\sum_{x_1\dots x_T}\frac{Q(\vec x)^2}{P_0(\vec x)}
=\E_{\psi,\psi'}\Big[\,G(\psi,\psi')^{T}\,\Big],\notag\\
G(\psi,\psi')&:=\sum_x\frac{p_\psi(x)\,p_{\psi'}(x)}{p_0(x)}
=d\sum_x\frac{\bra\psi E_x\ket\psi\,\bra{\psi'}E_x\ket{\psi'}}{\Tr E_x}.
\label{eq:chi2}
\end{align}
A direct estimate of \eqref{eq:chi2} over Haar pairs (the overlap $|\braket{\psi}{\psi'}|^2$ concentrates at $1/d$) yields $G=1+\bigO(1/d)$ off a measure-$\bigO(1/d)$ event, whence $\chi^2\le \exp(cT^2/d)-1$. This is the engine of the following theorem, whose tight, fully adaptive form is due to Chen--Cotler--Huang--Li and Huang et al.

\begin{theorem}[Incoherent purity-testing lower bound \cite{CCHL,LearnExp}]\label{thm:t1low}
Any incoherent (single-copy, adaptive) learner that distinguishes $\Tr(\rho^2)=1$ from $\Tr(\rho^2)\le\tfrac12$ with success probability $\ge 2/3$ requires
\begin{equation}
T\;\ge\;c\,2^{N/2}
\end{equation}
copies, for an absolute constant $c>0$. By contrast, coherent two-copy access succeeds with $T=\bigO(1)$.
\end{theorem}
\begin{proof}[Proof sketch]
Reduce to the two-point problem $H_0=\id/d$ vs.\ $H_1=$ Haar pure (Le Cam). By Lemma~\ref{lem:blind} the per-copy means match, so the likelihood ratio has no first-order term; the distinguishing information lives in \eqref{eq:chi2}. Bounding $\E_{\psi,\psi'}[G^T]$ via Lemma~\ref{lem:2design} gives $\chi^2(Q\Vert P_0)=\bigO(T^2/d)$, so $T=\Omega(\sqrt d)=\Omega(2^{N/2})$ is necessary for constant total-variation distance. The adaptive case is handled by the martingale/tree argument of \cite{CCHL}. The two-copy upper bound is Proposition~\ref{prop:t1up} with $\eps=\tfrac14$.\end{proof}

\subsection{Separation and noise robustness}
\begin{corollary}[Exponential sample separation]\label{cor:t1}
For purity testing on $N$ qubits, the platform uses $M_Q=\bigO(1)$ two-copy measurements while every quantum-memory-free strategy needs $\Omega(2^{N/2})$, an \emph{exponential} separation. For $\eps$-accurate estimation, $M_Q(\eps,\delta)=\lceil \ln(2/\delta)/(2\eps^2)\rceil$ is $N$-independent.
\end{corollary}

\begin{proposition}[Willow noise condition]\label{prop:t1noise}
The two-copy primitive uses $g_2=\bigO(N)$ native two-qubit gates and $r=\bigO(N)$ read-outs (Bell-basis variant), so by \eqref{eq:errbudget} the estimator bias is $\le\Gamma=\bigO\big((\eta_2+\eta_{\mathrm{ro}})N\big)$. The exponential separation of Corollary~\ref{cor:t1} therefore survives whenever
\begin{equation}
\eps\;\gtrsim\;(\eta_2+\eta_{\mathrm{ro}})\,N .
\label{eq:t1noise}
\end{equation}
With representative Willow numbers ($\eta_2+\eta_{\mathrm{ro}}\approx 10^{-2}$), the platform retains a dimension-independent $\eps\approx 0.05$ estimate up to $N\approx 5$ qubits per copy without error correction, and to arbitrary $N$ once the read-out is fault-tolerant.
\end{proposition}

\begin{remark}[Why the platform is the right vehicle]
The upper bound needs a stable second copy and repeated coherent access; Lemma~\ref{lem:sew} supplies persistent memory and Assumption~\ref{ass:idepth} keeps the Bell circuit at constant depth, so the $N$-independent cost of Proposition~\ref{prop:t1up} is physically attainable in a streaming setting rather than only in principle.
\end{remark}

%==========================================================================
\section{Task 2: quantum-enhanced sensing}\label{sec:task2}

\subsection{Estimation-theoretic model}
A field couples to $N$ probes through $U_\phi=e^{-i\phi G}$ with generator $G=\tfrac12\sum_{i=1}^N Z_i$. From $M$ repetitions one forms an estimator $\widehat\phi$. The benchmark is the quantum Cram\'er--Rao bound (QCRB)
\begin{equation}
\Var(\widehat\phi)\;\ge\;\frac{1}{M\,\FQ(\rho_\phi)},\qquad
\FQ=\Tr(\rho\,L^2),\quad \tfrac12(L\rho+\rho L)=\partial_\phi\rho_\phi,
\label{eq:qcrb}
\end{equation}
with $L$ the symmetric logarithmic derivative.

\begin{lemma}[Pure-state QFI]\label{lem:qfi}
For $\ket{\psi_\phi}=e^{-i\phi G}\ket{\psi_0}$, $\;\FQ=4\big(\bra{\partial_\phi\psi}\partial_\phi\psi\rangle-|\bra{\psi}\partial_\phi\psi\rangle|^2\big)=4\,\Var_{\psi_0}(G).$
\end{lemma}
\begin{proof}
$\ket{\partial_\phi\psi}=-iG\ket{\psi_\phi}$, so $\braket{\partial_\phi\psi}{\partial_\phi\psi}=\bra{\psi_\phi}G^2\ket{\psi_\phi}$ and $\bra\psi\partial_\phi\psi\rangle=-i\bra{\psi_\phi}G\ket{\psi_\phi}$. Substituting gives $4(\langle G^2\rangle-\langle G\rangle^2)$.\end{proof}

\subsection{Achievability: Heisenberg scaling on the platform}
\begin{proposition}[GHZ memory attains the Heisenberg limit]\label{prop:t2up}
Let the platform prepare the GHZ probe $\ket{\mathrm{GHZ}}=(\ket{0}^{\otimes N}+\ket{1}^{\otimes N})/\sqrt2$ via the native nearest-neighbour CZ/CNOT chain. Then $\FQ=N^2$, so $\Var(\widehat\phi)\ge 1/(MN^2)$, and the parity (Ramsey) read-out $\Pi=\prod_i X_i$ saturates the QCRB at small $\phi$: the signal $\langle\Pi\rangle=\cos(N\phi)$ has classical Fisher information $N^2$.
\end{proposition}
\begin{proof}
$(\,\sum_i Z_i)\ket{\mathrm{GHZ}}=\tfrac1{\sqrt2}(N\ket{0^N}-N\ket{1^N})$, so $\langle G\rangle=0$ and $\langle G^2\rangle=\tfrac14 N^2$, giving $\Var_{\mathrm{GHZ}}(G)=N^2/4$ and, by Lemma~\ref{lem:qfi}, $\FQ=N^2$. For the parity read-out, $\partial_\phi\langle\Pi\rangle=-N\sin(N\phi)$ and $\Var(\Pi)=1-\cos^2(N\phi)$, so the classical Fisher information $I(\phi)=(\partial_\phi\langle\Pi\rangle)^2/\Var(\Pi)=N^2$, matching $\FQ$.\end{proof}

\subsection{Lower bound without entangled memory}
\begin{theorem}[Standard quantum limit for separable probes \cite{GLM,Toth}]\label{thm:sql}
For any separable (unentangled) input state of the $N$ probes evolving under the local generator $G=\tfrac12\sum_i Z_i$, the QFI obeys $\FQ\le N$; hence every such strategy satisfies $\Var(\widehat\phi)\ge 1/(MN)$.
\end{theorem}
\begin{proof}[Proof sketch]
QFI is convex in the state and additive over a tensor product of independent probes; for a single qubit under $\tfrac12 Z$, $\FQ\le 1$ with equality on an equatorial state. Convexity extends the product bound $\FQ\le N$ to all separable states.\end{proof}

\begin{corollary}[Quadratic metrological separation]\label{cor:t2}
The platform attains $\Var(\widehat\phi)=\Theta(1/(MN^2))$; every quantum-memory-free (separable) scheme is bounded by $\Var(\widehat\phi)\ge 1/(MN)$. Equivalently, to reach a target variance $\sigma^2$ the platform needs a factor $N$ fewer probe interrogations.
\end{corollary}

\subsection{Beyond QFI: the Grover--Heisenberg ceiling, and noise}
For \emph{unknown-frequency} detection of a weak oscillating field, processing the digitised signal on the quantum computer is strictly stronger than any quantum-Fisher-information or finite-lifetime-memory protocol; the optimal sensing time is set by the Grover--Heisenberg limit \cite{QCES}. The platform's coherent memory is exactly the resource that paper's lower bound presumes, so Corollary~\ref{cor:t2} is the conservative (fixed-frequency) face of a larger advantage.

\begin{proposition}[Dephasing trade-off]\label{prop:t2noise}
Under independent dephasing at rate $\gamma$ during interrogation time $t$, the GHZ coherence decays as $e^{-N\gamma t}$, giving effective QFI $\FQ^{\mathrm{eff}}(t)=N^2 t^2 e^{-2N\gamma t}$, maximised at $t^\star=1/(N\gamma)$ with $\FQ^{\mathrm{eff}}(t^\star)=N/(e^2\gamma^2)$. Thus under uncorrected Markovian dephasing the advantage degrades to a \emph{constant-factor} improvement \cite{Huelga}; under Willow-class error correction (logical dephasing $\gamma_L\ll\gamma$, achievable below threshold \cite{Willow}) the full $1/N^2$ scaling is restored up to the logical coherence time. The crossover qubit number is $N^\star\approx (\gamma t_{\mathrm{gate}})^{-1}$.
\end{proposition}

%==========================================================================
\section{Task 3: tracking entanglement dynamics}\label{sec:task3}

\subsection{Model}
A many-body state $\rho(t)=U_t\rho(0)U_t^\dagger$ evolves; the target is the subsystem second R\'enyi entropy of a region $A$ of $|A|$ qubits,
\begin{equation}
S_2^A(t)=-\log_2 P_A(t),\qquad P_A(t)=\Tr\!\big(\rho_A(t)^2\big),\quad \rho_A=\Tr_{\bar A}\rho .
\end{equation}
$P_A$ is a two-copy functional, so Task~3 is the temporal, subsystem-restricted instance of Task~1.

\subsection{Achievability}
\begin{proposition}[Per-step, dimension-independent tracking]\label{prop:t3up}
At each time $t$, the platform holds $\rho_A(t)$ on two copies and applies the SWAP test restricted to $A$ (Lemma~\ref{lem:swap}), giving an unbiased $\widehat P_A(t)$ with $|\widehat P_A-P_A|\le\eps$ at confidence $1-\delta'$ using $M_Q=\lceil\ln(2/\delta')/(2\eps^2)\rceil$ shots, \emph{independent of $|A|$ and $N$}. Over a horizon of $T$ steps, a union bound with $\delta'=\delta/T$ gives uniform additive accuracy $\eps$ with $M_{\mathrm{tot}}=\bigO\!\big(\eps^{-2}T\log(T/\delta)\big)$ total shots. If $P_A(t)\ge p_{\min}$ on the horizon, then $|\widehat S_2^A-S_2^A|\le \eps/(p_{\min}\ln 2)+\bigO(\eps^2)$.
\end{proposition}
\begin{proof}
The per-step bound is Proposition~\ref{prop:t1up} applied to the $|A|$-qubit subsystem (the SWAP acts only on $A$). The union bound controls all $T$ steps simultaneously. The entropy statement is first-order propagation of error through $S_2=-\log_2 P_A$, whose derivative magnitude is $1/(P_A\ln2)\le 1/(p_{\min}\ln2)$.\end{proof}

\subsection{Lower bound without quantum memory}
\begin{theorem}[Incoherent subsystem-purity lower bound]\label{thm:t3low}
Estimating $P_A=\Tr(\rho_A^2)$ to additive constant error with incoherent (single-copy) access requires $\Omega(2^{|A|/2})$ copies in the worst case over states on $A$; consequently any quantum-memory-free tracker needs $\Omega(T\,2^{|A|/2})$ copies over a $T$-step horizon.
\end{theorem}
\begin{proof}
Apply Theorem~\ref{thm:t1low} with $N\to|A|$: the hard ensemble can be embedded in $A$ while $\bar A$ is held in a fixed product state, so the subsystem-purity problem inherits the $\Omega(2^{|A|/2})$ incoherent lower bound. The horizon factor is the $T$ independent steps.\end{proof}

\begin{corollary}[Exponential-in-$|A|$ separation, uniform over the trajectory]\label{cor:t3}
Per step, the platform's cost $M_Q=\bigO(\eps^{-2})$ is independent of $|A|$, whereas every quantum-memory-free estimator costs $\Omega(2^{|A|/2})$. The classical-shadow blow-up observed numerically (tracking RMSE growing from $\sim0.18$ at $|A|=1$ to $\sim5.6$ at $|A|=3$ at fixed shots) is the finite-size signature of this theorem.
\end{corollary}
\begin{proposition}[Willow noise condition]
The subsystem SWAP/Bell circuit uses $g_2=\bigO(|A|)$ native gates, so the per-step bias is $\Gamma=\bigO((\eta_2+\eta_{\mathrm{ro}})|A|)$ and the separation of Corollary~\ref{cor:t3} survives for $\eps\gtrsim(\eta_2+\eta_{\mathrm{ro}})|A|$; entropies remain accurate while $P_A\ge p_{\min}\gg\Gamma$.
\end{proposition}

%==========================================================================
\section{Task 4: real-time QEC decoding}\label{sec:task4}

The inputs (syndromes) are classical bitstrings, so Task~4 carries \emph{no} exponential quantum-data separation; we state precisely the guarantees it \emph{does} admit: (i) a representational/learning guarantee for the reservoir read-out, and (ii) an unconditional real-time (no-backlog) latency guarantee.

\subsection{Learnability of the optimal finite-history decoder}
\begin{definition}
A decoder is a causal map from the syndrome stream $\{s_\tau\}_{\tau\le t}$ to a logical correction. A \emph{fading-memory} decoder is one whose dependence on $s_{t-k}$ decays with $k$; the maximum-likelihood decoder restricted to a window $W$ is fading-memory.
\end{definition}

\begin{theorem}[Reservoir universality for fading-memory decoders \cite{GO,BoydChua}]\label{thm:t4univ}
A reservoir with the echo-state (fading-memory) property, read out by a linear map on a polynomially rich feature set, is a universal approximator of fading-memory filters: for any such target decoder $f^\star$ and any $\eps>0$ there exist read-out weights $w$ with $\sup\|w\!\cdot\!\phi(\cdot)-f^\star\|\le\eps$. The critical-phase reservoir \eqref{eq:Ustep} at $g^\star$ has the echo-state property and, being barren-plateau-free under local read-out, admits a trainable such $w$.
\end{theorem}

\begin{proposition}[Generalisation of the trained read-out]\label{prop:t4gen}
With $F$ reservoir features bounded by $\|\phi\|\le B$, labels in $\{0,1\}$, and $M$ i.i.d.\ training sequences, ridge/logistic read-out has excess risk
\begin{equation}
\E[\mathcal R(\widehat w)]-\min_w \mathcal R(w)\;\le\;\bigO\!\Big(B\sqrt{F/M}\Big),
\end{equation}
by the Rademacher complexity of linear predictors. Only $F$ read-out weights are trained --- no backpropagation through a deep network, in contrast to a recurrent-transformer decoder.
\end{proposition}

\subsection{The real-time (no-backlog) guarantee}
\begin{theorem}[Backlog stability]\label{thm:backlog}
Let $\tau_{\mathrm{gen}}$ be the QEC cycle period (one syndrome frame per cycle) and $\tau_{\mathrm{dec}}$ the per-frame decode latency. The decoding backlog $b_k$ obeys $b_{k+1}=\max\{0,\,b_k+\tau_{\mathrm{dec}}-\tau_{\mathrm{gen}}\}$ and remains bounded (real-time decoding is sustainable) \emph{iff} $\tau_{\mathrm{dec}}\le\tau_{\mathrm{gen}}$. If $\tau_{\mathrm{dec}}>\tau_{\mathrm{gen}}$ the backlog diverges linearly, $b_k\sim k(\tau_{\mathrm{dec}}-\tau_{\mathrm{gen}})$, the ``exponential stalling'' regime.
\end{theorem}
\begin{proof}
The recurrence is the Lindley equation of a deterministic single-server queue with inter-arrival time $\tau_{\mathrm{gen}}$ and service time $\tau_{\mathrm{dec}}$; its stationary backlog is finite iff utilisation $\tau_{\mathrm{dec}}/\tau_{\mathrm{gen}}\le 1$, and otherwise the unfinished work grows at rate $\tau_{\mathrm{dec}}-\tau_{\mathrm{gen}}$.\end{proof}

\begin{corollary}[Constant-latency decoding across code distance]\label{cor:t4}
The reservoir read-out is a single matrix--vector product of cost $\bigO(F)$, and the on-chip reservoir update has $\bigO(1)$ physical depth (Assumption~\ref{ass:idepth}); hence $\tau_{\mathrm{dec}}=\bigO(F)$ is \emph{independent of code distance $d$}. Therefore there exists a fixed feature budget $F$ with $\tau_{\mathrm{dec}}\le\tau_{\mathrm{gen}}$ for all $d$, so Theorem~\ref{thm:backlog} guarantees backlog-free decoding at every distance --- whereas any decoder whose latency grows with $d$ eventually violates stability.
\end{corollary}
\begin{remark}[Honest scope]
Corollary~\ref{cor:t4} bounds \emph{latency and training cost}, not the logical error rate; matching the accuracy of optimal matching/transformer decoders is empirical (the experiments show parity with a trained linear decoder and a large margin over trivial), not guaranteed here.
\end{remark}

%==========================================================================
\section{Task 5: generative modelling of a classically-hard process}\label{sec:task5}

We give two guarantees: an \emph{unconditional} memory separation and a \emph{conditional} sampling-hardness separation.

\subsection{Unconditional memory advantage}
\begin{definition}[Statistical complexity]
For a stationary process, the classical causal states are the equivalence classes of pasts with identical future conditionals; the minimal classical predictive memory is the statistical complexity $C_\mu=H(\mathcal S)$ (Shannon entropy of the causal states, the provably minimal classical model --- the $\eps$-machine). The minimal \emph{quantum} predictive memory is $C_q=S(\rho_{\mathcal S})\le C_\mu$, the von Neumann entropy of the quantum causal states.
\end{definition}

\begin{theorem}[Quantum models are more memory-efficient \cite{Gu,Garner}]\label{thm:t5mem}
There exist stationary processes with $C_q<C_\mu$, and parametrised families for which the ratio is unbounded, $C_\mu/C_q\to\infty$. Consequently a quantum predictive model can match the classical process's future statistics exactly while storing strictly --- and unboundedly --- less memory than the minimal classical model.
\end{theorem}

\begin{corollary}[Platform realisation]\label{cor:t5mem}
The reservoir's internal quantum state plays the role of the quantum causal states; the collapse-free read-out (Lemma~\ref{lem:sew}) maintains them across the stream without the projective reset that a classical model would require. Hence on such processes the platform predicts to the same accuracy with memory $C_q$, beating the classical lower bound $C_\mu$ by an unbounded factor.
\end{corollary}

\subsection{Conditional sampling-hardness}
The platform emits symbols by measuring $U_D\ket0$ for the depth-$D$ critical-phase circuit \eqref{eq:Ustep}; let $p(x)=|\bra x U_D\ket0|^2$ be the per-step output law.
\begin{assumption}[Anticoncentration + average-case hardness \cite{IQP,BFNV}]\label{ass:hard}
At depth $D\ge D_0$ the law $p$ anticoncentrates ($\Pr_x[p(x)\ge 1/(2d)]=\Omega(1)$) and computing $|\bra x U_D\ket0|^2$ to small relative error is $\#\mathrm P$-hard on average.
\end{assumption}
\begin{theorem}[No efficient classical next-symbol model]\label{thm:t5hard}
Under Assumption~\ref{ass:hard} and the non-collapse of the polynomial hierarchy, no classical $\poly(N)$-time algorithm samples the platform's per-step law $p$ to total-variation error $<1/\poly(N)$. Hence no efficient classical generative model can reproduce the platform's next-symbol conditionals to inverse-polynomial accuracy.
\end{theorem}
\begin{proof}[Proof sketch]
A classical sampler of $p$ to small TV, combined with anticoncentration (Assumption~\ref{ass:hard}) and Stockmeyer counting, would place an average-case $\#\mathrm P$-hard quantity inside $\mathrm{BPP}^{\mathrm{NP}}$, collapsing the polynomial hierarchy to its third level --- the standard IQP/random-circuit-sampling argument \cite{IQP,BFNV}. A classical model matching the conditionals to inverse-poly TV would furnish such a sampler, a contradiction. This is the sequential face of the generative advantage of \cite{GenAdv}.\end{proof}

%==========================================================================
\section{Summary of guarantees}\label{sec:summary}

\begin{center}
\renewcommand{\arraystretch}{1.35}
\begin{tabular}{@{}p{2.3cm}p{2.5cm}p{2.9cm}p{3.0cm}p{2.4cm}@{}}
\hline
\textbf{Task} & \textbf{Platform cost} & \textbf{Memory-free cost} & \textbf{Separation} & \textbf{Type}\\
\hline
1. Functionals & $\bigO(\eps^{-2}\log\tfrac1\delta)$, $N$-indep. & $\Omega(2^{N/2})$ & exponential in $N$ & unconditional\\
2. Sensing & $\Var=\Theta(1/MN^2)$ & $\Var\ge 1/MN$ & quadratic in $N$ & unconditional\\
3. Entanglement & $\bigO(\eps^{-2})$ per step, $|A|$-indep. & $\Omega(2^{|A|/2})$ per step & exp.\ in $|A|$, $\forall t$ & unconditional\\
4. Decoding & $\tau_{\mathrm{dec}}=\bigO(F)$, $d$-indep. & latency $\nearrow d\Rightarrow$ backlog & constant-latency, cheap training & unconditional (latency)\\
5. Generative & memory $C_q$; samples $p$ natively & memory $C_\mu$; no efficient sampler & unbounded memory ratio; sampling hardness & uncond.\ (memory) / cond.\ (sampling)\\
\hline
\end{tabular}
\end{center}

Every separation is realised by the same three platform primitives: the collapse-free read-out (persistent quantum memory, Lemma~\ref{lem:sew}), instantaneous depth (constant-depth circuits satisfying the noise budgets, Assumption~\ref{ass:idepth}), and local-cost trainability (the read-out in Tasks~4--5 is learnable). The unifying principle is the with-/without-quantum-memory dichotomy: Tasks~1 and~3 are its purest expression (second-order/non-linear functionals invisible to first-order, single-copy statistics), Task~2 its metrological instance (entangled vs.\ separable Fisher information), and Tasks~4--5 its practical (latency) and generative (memory/sampling) instances.

\section{Limitations and honesty}\label{sec:limits}
\begin{itemize}[leftmargin=1.4em]
\item Tasks~1, 3 require faithful loading of the quantum data into the memory; the separations are over \emph{sample/copy} complexity, assuming such access (as in \cite{LearnExp}).
\item Task~2's full $1/N^2$ needs coherence beyond uncorrected Willow dephasing (Proposition~\ref{prop:t2noise}); below threshold it is restored, otherwise a constant-factor gain remains.
\item Task~4 guarantees latency and learnability, \emph{not} optimal logical error rate; accuracy parity is empirical.
\item Task~5's sampling-hardness is conditional on standard complexity conjectures (Assumption~\ref{ass:hard}, PH non-collapse); the memory advantage is unconditional.
\item All proofs are stated for the idealised platform; finite Willow error enters only through the budget $\Gamma$ of \eqref{eq:errbudget}, and each separation is paired with the explicit condition under which $\Gamma$ is tolerable.
\end{itemize}

\begin{thebibliography}{99}
\bibitem{LearnExp} H.-Y.\ Huang \emph{et al.}, ``Quantum advantage in learning from experiments,'' \emph{Science} \textbf{376}, 1182 (2022); arXiv:2112.00778.
\bibitem{CCHL} S.\ Chen, J.\ Cotler, H.-Y.\ Huang, J.\ Li, ``Exponential separations between learning with and without quantum memory,'' \emph{FOCS} 2021; arXiv:2111.05881.
\bibitem{LSQC} H.-Y.\ Huang \emph{et al.}, ``Learning shallow quantum circuits,'' \emph{STOC} 2024; arXiv:2401.10095.
\bibitem{GenAdv} Huang \emph{et al.}, ``Generative quantum advantage for classical and quantum problems'' (sewing; instantaneous depth), arXiv:2509.09033 (2025).
\bibitem{HKP} H.-Y.\ Huang, R.\ Kueng, J.\ Preskill, ``Predicting many properties of a quantum system from very few measurements,'' \emph{Nature Physics} \textbf{16}, 1050 (2020).
\bibitem{GLM} V.\ Giovannetti, S.\ Lloyd, L.\ Maccone, ``Advances in quantum metrology,'' \emph{Nature Photonics} \textbf{5}, 222 (2011).
\bibitem{Toth} G.\ T\'oth, I.\ Apellaniz, ``Quantum metrology from a quantum information science perspective,'' \emph{J.\ Phys.\ A} \textbf{47}, 424006 (2014).
\bibitem{QCES} R.\ R.\ Allen, F.\ Machado, I.\ L.\ Chuang, H.-Y.\ Huang, S.\ Choi, ``Quantum computing enhanced sensing,'' arXiv:2501.07625 (2025).
\bibitem{Huelga} S.\ F.\ Huelga \emph{et al.}, ``Improvement of frequency standards with quantum entanglement,'' \emph{Phys.\ Rev.\ Lett.} \textbf{79}, 3865 (1997).
\bibitem{Willow} Google Quantum AI, ``Quantum error correction below the surface-code threshold'' (Willow), \emph{Nature} \textbf{638}, 920 (2025).
\bibitem{GO} L.\ Grigoryeva, J.-P.\ Ortega, ``Echo state networks are universal,'' \emph{Neural Networks} \textbf{108}, 495 (2018).
\bibitem{BoydChua} S.\ Boyd, L.\ Chua, ``Fading memory and the problem of approximating operators with Volterra series,'' \emph{IEEE Trans.\ Circuits Syst.} \textbf{32}, 1150 (1985).
\bibitem{Gu} M.\ Gu, K.\ Wiesner, E.\ Rieper, V.\ Vedral, ``Quantum mechanics can reduce the complexity of classical models,'' \emph{Nature Communications} \textbf{3}, 762 (2012).
\bibitem{Garner} A.\ J.\ P.\ Garner \emph{et al.}, ``Provably unbounded memory advantage in stochastic simulation using quantum mechanics,'' \emph{New J.\ Phys.} \textbf{19}, 103009 (2017).
\bibitem{IQP} M.\ J.\ Bremner, A.\ Montanaro, D.\ J.\ Shepherd, ``Average-case complexity versus approximate simulation of commuting quantum computations,'' \emph{Phys.\ Rev.\ Lett.} \textbf{117}, 080501 (2016).
\bibitem{BFNV} A.\ Bouland, B.\ Fefferman, C.\ Nirkhe, U.\ Vazirani, ``On the complexity and verification of quantum random circuit sampling,'' \emph{Nature Physics} \textbf{15}, 159 (2019).
\end{thebibliography}

\end{document}
